% TEMPLATE-MILC-v3.tex - plantilla maestra para todas las guias MILC v3
% Ver scripts/generadores/build-guia-milc.py para la lista completa de
% claves esperadas. El script valida que no queden placeholders sin
% reemplazar antes de escribir el archivo final.

\documentclass[11pt,letterpaper]{article}
\usepackage[margin=1.75cm]{geometry}
\usepackage[table]{xcolor}
\usepackage{fontspec}
\usepackage[spanish]{babel}
\usepackage{tikz}
\usepackage[most]{tcolorbox}
\usepackage{tabularx}
\usepackage{array}
\usepackage{fancyhdr}
\usepackage{titlesec}
\usepackage{ragged2e}
\usepackage{enumitem}
\usepackage{microtype}

\setmainfont{Helvetica Neue}
\setsansfont{Helvetica Neue}
\setlength{\parindent}{0pt}
\setlength{\parskip}{6pt}
\hyphenpenalty=200
\exhyphenpenalty=200
\pretolerance=400
\tolerance=2000
\setlength{\emergencystretch}{1em}
\renewcommand{\arraystretch}{1.25}
\setlist{leftmargin=5mm,itemsep=1mm,topsep=1mm}
\newcolumntype{Y}{>{\RaggedRight\arraybackslash}X}

\definecolor{milcMagenta}{HTML}{E6007E}
\definecolor{milcVino}{HTML}{5A0038}
\definecolor{milcTurquesa}{HTML}{0093A5}
\definecolor{milcVerde}{HTML}{58A923}
\definecolor{milcMostaza}{HTML}{E5B400}
\definecolor{milcOcre}{HTML}{B86600}
\definecolor{milcNegro}{HTML}{241816}
\definecolor{milcGris}{HTML}{F7F7F4}
\definecolor{milcLinea}{HTML}{E8E2DD}
\definecolor{milcGradePrimary}{HTML}{84CC16}
\definecolor{milcGradeDark}{HTML}{4D7C0F}
\definecolor{milcGradeSoft}{HTML}{ECFCCB}
\definecolor{milcGradeAccent}{HTML}{CFE7FF}
\definecolor{milcGradeText}{HTML}{FFFFFF}
\color{milcNegro}

\pagestyle{fancy}
\fancyhf{}
\lhead{\small Tecnología e Informática · Bebras}
\rhead{\small Momento 6 · MILC v3}
\cfoot{\small\thepage}
\renewcommand{\headrulewidth}{0pt}

\newtcolorbox{titlebox}[1]{
  enhanced,colback=#1,colframe=#1,arc=7mm,boxrule=0pt,
  left=7mm,right=7mm,top=3mm,bottom=3mm,
  before skip=8pt,after skip=8pt,
  fontupper=\sffamily\bfseries\Large\color{white}
}

\newtcolorbox{softbox}[3]{
  enhanced,colback=#2,colframe=#1,borderline west={4pt}{0pt}{#1},
  arc=2mm,boxrule=.4pt,left=4mm,right=4mm,top=3mm,bottom=3mm,
  before skip=6pt,after skip=8pt,title={#3},coltitle=white,
  colbacktitle=#1,fonttitle=\sffamily\bfseries,fontupper=\RaggedRight,
  attach boxed title to top left={xshift=3mm,yshift=-2mm},
  boxed title style={arc=2mm, boxrule=0pt}
}

\newtcolorbox{drawbox}[2]{
  enhanced,colback=white,colframe=#1,arc=4mm,boxrule=1pt,
  left=5mm,right=5mm,top=4mm,bottom=4mm,before skip=8pt,after skip=8pt,
  title={#2},coltitle=white,colbacktitle=#1,fonttitle=\sffamily\bfseries,
  fontupper=\RaggedRight,
  attach boxed title to top left={xshift=4mm,yshift=-2mm},
  boxed title style={arc=3mm, boxrule=0pt}
}

\newtcolorbox{infoband}[1]{
  enhanced,colback=#1!8,colframe=#1!50,arc=2mm,boxrule=.5pt,
  left=4mm,right=4mm,top=2mm,bottom=2mm,before skip=4pt,after skip=4pt,
  fontupper=\small\RaggedRight
}

% Puente narrativo entre secciones (contrato editorial Conéctate).
% Da continuidad: "Acabas de X. Ahora vas a Y porque Z."
% Visualmente: banda izquierda en color del grado, texto en itálica pequeña.
\newtcolorbox{puentebox}{
  enhanced,colback=milcGradeSoft,colframe=milcGradeDark,
  borderline west={3pt}{0pt}{milcGradeDark},
  arc=0mm,boxrule=0pt,left=5mm,right=4mm,top=2.5mm,bottom=2.5mm,
  before skip=4pt,after skip=8pt,
  fontupper=\itshape\small\color{milcNegro}\RaggedRight
}
\newcommand{\puente}[1]{%
  \begin{puentebox}%
    \textcolor{milcGradeDark}{$\rightarrow$}\hspace{2mm}#1%
  \end{puentebox}%
}

\newcommand{\linead}[1]{\noindent\color{milcLinea}\rule{#1}{0.5pt}\color{milcNegro}}
\newcommand{\respuesta}[1]{\par\vspace{2mm}\foreach \n in {1,...,#1}{\noindent\color{milcLinea}\rule{\linewidth}{0.5pt}\par\vspace{5mm}}\color{milcNegro}}
\newcommand{\checkbox}{\textcolor{milcGradeDark}{\fbox{\phantom{x}}}\hspace{5pt}}

\newcommand{\phasecard}[4]{
\begin{tcolorbox}[enhanced,colback=#3,colframe=#2,arc=3mm,boxrule=.5pt,height=3.05cm,valign=center,halign=center,left=1.5mm,right=1.5mm,top=2mm,bottom=2mm]
{\sffamily\bfseries\fontsize{22}{24}\selectfont\textcolor{#2}{#1}}\par
{\sffamily\bfseries\small #4}
\end{tcolorbox}
}

\begin{document}
\thispagestyle{empty}
\begin{tikzpicture}[remember picture,overlay]
  \fill[white] (current page.south west) rectangle (current page.north east);
  \path[fill=milcGradePrimary,rounded corners=16mm]
    ([xshift=.55cm,yshift=-.55cm]current page.north west) rectangle
    ([xshift=-.55cm,yshift=3.65cm]current page.south east);

  \node[anchor=north west,text=milcGradeText] at ([xshift=2.0cm,yshift=-1.85cm]current page.north west)
    {\sffamily\bfseries\fontsize{14}{16}\selectfont MOMENTO};
  \node[anchor=north west,text=milcGradeText,opacity=.82] at ([xshift=2.0cm,yshift=-2.75cm]current page.north west)
    {\sffamily\bfseries\fontsize{9}{11}\selectfont BEBRAS · PENSAR COMO UNA MÁQUINA};

  \begin{scope}[shift={([xshift=-3.05cm,yshift=-2.55cm]current page.north east)}]
    \fill[white,opacity=.80] (-.62,-.52) rectangle (.62,.52);
    \draw[milcGradeText,line width=1pt] (-.62,-.52) rectangle (.62,.52);
    \fill[milcGradeDark] (-.42,-.38) rectangle (-.22,.06);
    \fill[milcGradeAccent] (-.10,-.38) rectangle (.10,.24);
    \fill[milcGradePrimary!60!black] (.22,-.38) rectangle (.42,.38);
  \end{scope}

  \node[anchor=west,text=milcGradeText] at ([xshift=1.9cm,yshift=-12.85cm]current page.north west)
    {\sffamily\bfseries\fontsize{124}{124}\selectfont 6};
  \draw[milcGradeText,line width=1.7pt]
    ([xshift=4.52cm,yshift=-11.68cm]current page.north west) circle (.13cm);

  \node[anchor=west,text=milcGradeText,text width=8.7cm,align=left] at ([xshift=10.35cm,yshift=-11.6cm]current page.north west)
    {
     {\sffamily\bfseries\fontsize{19}{22}\selectfont Verdadero o falso --- lógica,\\deducción y acertijos de grilla}\\[4mm]
     {\sffamily\bfseries\fontsize{23}{26}\selectfont Bebras · Momento 6}\\[2mm]
     {\sffamily\fontsize{13}{16}\selectfont Curso «Pensar como una máquina» · Pensamiento computacional}\\[5mm]
     {\sffamily\bfseries\fontsize{11}{13}\selectfont Institución Educativa Sor María Juliana}\\[1mm]
     {\sffamily\fontsize{10}{12}\selectfont Autor: Dr. Álvaro Cárdenas Orozco}
    };

  \path[fill=milcGradeSoft,rounded corners=7mm]
    ([xshift=1.35cm,yshift=.85cm]current page.south west) rectangle
    ([xshift=-1.35cm,yshift=4.25cm]current page.south east);
  \draw[milcGradeDark,line width=.65pt,rounded corners=7mm]
    ([xshift=1.35cm,yshift=.85cm]current page.south west) rectangle
    ([xshift=-1.35cm,yshift=4.25cm]current page.south east);
  \node[anchor=center,text=milcNegro,text width=17.0cm,align=center] at ([yshift=2.55cm]current page.south)
    {\sffamily\fontsize{10}{12}\selectfont Metodología MILC · Apertura ancestral · Pensamiento computacional · Triángulo Dussel-Estoicismo-Floridi\\
    \textcolor{milcGradeDark}{\bfseries Producto:} El acertijo de los tres barrios resuelto con grilla 3×3 y el veredicto escrito sobre una deducción inválida, con justificación paso a paso.};
\end{tikzpicture}
\mbox{}
\newpage

\section*{Apertura: Verdadero o falso: lógica, deducción y acertijos de grilla}

\begin{softbox}{milcGradeDark}{milcGradeSoft}{Saber ancestral}
En los patios de Cartago, en las cocinas a leña del Valle del Cauca y en las noches del Pacífico colombiano, los abuelos no enseñaban a pensar con cuadernos: enseñaban con \textbf{adivinanzas}. «Agua pasa por mi casa, cate de mi corazón»: quien la escucha no adivina al azar, \emph{deduce}. Tiene pistas --- «agua» al comienzo, «cate» al final --- y, juntándolas con lógica, llega a la única respuesta que las cumple todas: el aguacate. Una buena adivinanza nunca tiene dos respuestas: da exactamente las señas necesarias para descartar lo imposible hasta que solo queda lo verdadero. Así los mayores afilaban la mente de los niños sin que pareciera estudio. El acertijo de «quién tiene qué», el de las tres puertas, el del que miente y el que dice verdad: todos piden lo mismo, descartar lo falso, juntar lo que sabes y quedarte con la única salida posible. \textbf{Origen:} tradición oral afrocolombiana y campesina del Valle del Cauca y el Pacífico colombiano.
\end{softbox}

\begin{softbox}{milcTurquesa}{milcGris}{Saber contemporáneo}
\textbf{Lógica} es razonar con dos valores --- verdadero y falso --- para deducir con certeza. No es opinión ni intuición: si las pistas son ciertas y razonas bien, la conclusión \emph{tiene que} ser cierta. La pieza central es el \textbf{condicional}: «SI llueve ENTONCES llevo paraguas». Esa regla solo se lee hacia adelante. Si sabes que \textbf{NO} llevas paraguas, puedes deducir que NO llovió --- eso se llama \textbf{contrapositiva} ---. Pero de «no llueve» no se deduce si llevas paraguas o no. Un \textbf{acertijo de grilla} aplica esta misma lógica: dibujás una tabla con personas en las filas y objetos en las columnas, marcás con X lo que cada pista descarta y, cuando solo queda una posibilidad por fila, ya tienes la respuesta --- con certeza, no por suerte ---. Muchas tareas Bebras son exactamente esto.
\end{softbox}

\begin{softbox}{milcMostaza}{milcGris}{Pregunta-puente}
\textit{Cuando deducís el aguacate de una adivinanza, descartás lo falso hasta que solo queda lo verdadero. ¿Y si te dijera que eso --- separar lo verdadero de lo falso para deducir con certeza --- es exactamente lo que hace por dentro un computador, y lo que te van a pedir muchas tareas Bebras?}
\end{softbox}

\begin{softbox}{milcVerde}{milcGris}{Saber y saber hacer}
Al terminar podrás: (1) \textbf{identificar} las cuatro herramientas de la lógica --- condicional, negación, Y/O y deducción por descarte ---; (2) \textbf{analizar} un acertijo de grilla marcando con X cada descarte que una pista permite, hasta llegar a la única asignación posible; (3) \textbf{evaluar} si una conclusión está bien deducida o si lee el condicional al revés, y justificar por escrito el veredicto; (4) aplicar la contrapositiva para resolver retos de nivel C en el simulacro estilo Bebras.
\end{softbox}

\small
\begin{tabularx}{\textwidth}{>{\bfseries}p{3.0cm}Y}
\rowcolor{milcGradePrimary!12}Autor & Dr. Álvaro Cárdenas Orozco\\
DBA & Desarrolla el pensamiento computacional --- descomposición, patrones, abstracción y algoritmos --- para resolver problemas (transversal 6°--11°, alineado a los lineamientos MEN de Tecnología e Informática)\\
Referentes & Valentina Dagienė (Bebras) · Pensamiento computacional (Wing) · Dussel · Estoicismo · Floridi · Saberes del Valle y el Pacífico\\
Producto final & El acertijo de los tres barrios resuelto con grilla 3×3 y el veredicto escrito sobre una deducción inválida, con justificación paso a paso.\\
\end{tabularx}
\normalsize

\puente{Ya sabes que la adivinanza afina la mente descartando lo falso. Ahora mira el mapa de la sesión: vas a recorrer las herramientas de la lógica, una por una, hasta resolver un acertijo completo de grilla.}

\begin{titlebox}{milcGradeDark}
Ruta MILC: así vamos a aprender
\end{titlebox}
\begin{tabularx}{\textwidth}{>{\centering\arraybackslash}X>{\centering\arraybackslash}X>{\centering\arraybackslash}X>{\centering\arraybackslash}X}
\phasecard{1}{milcVerde}{milcGris}{Escucha\\[-1mm]\small Observo, escucho y pregunto.} &
\phasecard{2}{milcTurquesa}{milcGris}{Sistematización\\[-1mm]\small Pensamiento computacional.} &
\phasecard{3}{milcMagenta}{milcGris}{Praxis\\[-1mm]\small Construyo y aplico.} &
\phasecard{4}{milcVino}{milcGris}{Evaluación\\[-1mm]\small Reflexión liberadora.}
\end{tabularx}


\puente{Antes de nombrar las herramientas, conviene verlas en acción. Empieza por escuchar una adivinanza clásica del Valle y preguntarte cómo la resolviste --- no solo qué respondiste.}

\begin{titlebox}{milcVerde}
Fase 1 · Escucha
\end{titlebox}

\begin{softbox}{milcVerde}{milcGris}{Escena para escuchar}
\textbf{Escuchá esta adivinanza} de la tradición oral colombiana: «Tengo dientes y no como, peino a los demás y a mí nadie me peina». Antes de responder, anda por partes: ¿qué descartás porque no cumple la primera pista? ¿Qué descartás porque no cumple la segunda? ¿Qué opción sobrevive a las dos? Fíjate que no adivinas al azar: \emph{deducís}. Después recordá o pídele a alguien cercano una adivinanza del Valle o del Pacífico, y resolvela con el mismo método --- descartando, no adivinando.
\end{softbox}

\checkbox Respondiste la adivinanza del peine y escribiste qué pista usaste para descartar cada opción que no funcionaba.\\
\checkbox Identificaste en tu adivinanza familiar al menos dos pistas y explicaste en qué orden las usaste para llegar a la respuesta.\\
\checkbox Nombraste, con tus palabras, la diferencia entre adivinar al azar y deducir con pistas.

\begin{infoband}{milcVerde}
\textbf{En tu cuaderno (Actividad 1 · IDENTIFICA).} Título: «Actividad 1 --- La adivinanza que deduzco». \textbf{Formato:} lista de 4 a 6 viñetas. \textbf{Extensión:} una viñeta por pista usada y lo que descartaste con ella. \textbf{Sabes que terminaste cuando} tu lista nombra, en lenguaje cotidiano, al menos dos descartes y concluye con la única respuesta que sobrevivió.
\end{infoband}

\puente{Acabás de ver la lógica en lo cotidiano. Ahora vamos a nombrar sus cuatro herramientas y entenderlas una a una, porque con ellas vas a resolver todo lo que sigue.}

\begin{titlebox}{milcTurquesa}
Fase 2 · Sistematización con pensamiento computacional
\end{titlebox}

\begin{softbox}{milcTurquesa}{milcGris}{Cuatro pilares aplicados al tema}
Deducir como en una adivinanza --- o como en una tarea Bebras --- se hace con cuatro herramientas. No es magia ni requiere código: son piezas que, juntas, convierten un acertijo confuso en algo que sí se puede resolver. Acá las nombramos una a una, con ejemplos de tu barrio en Cartago.
\end{softbox}

\small
\begin{tabularx}{\textwidth}{>{\bfseries\RaggedRight\arraybackslash}p{3.6cm}Y}
\rowcolor{milcTurquesa!90}\color{white}Pilar & \color{white}\bfseries Aplicación al tema\\
1. Descomposición & \textbf{Condicional (SI\ldots ENTONCES\ldots):} una regla que ata dos cosas. «SI llueve ENTONCES llevo paraguas.» No dice que siempre llueva: dice que \emph{cada vez} que llueve, hay paraguas. Es la pieza con la que se deduce. Cuidado: la regla solo va hacia adelante --- de «hay paraguas» no puedes concluir que llovió.\\
2. Reconocimiento de patrones & \textbf{Negación (NO):} voltea el valor de verdad. Si «está lloviendo» es verdadero, «NO está lloviendo» es falso. Saber que algo NO ocurre suele ser tan útil como saber que sí ocurre. Combinada con el condicional, da la \textbf{contrapositiva}: si no pasa lo segundo, tampoco pasa lo primero.\\
3. Abstracción & \textbf{Y / O:} «Y» pide que se cumplan las dos cosas a la vez --- tener arroz Y tener pollo ---. «O» pide al menos una de las dos --- pagar en efectivo O por Nequi ---. Cambian por completo lo que puedes concluir de una pista.\\
4. Algoritmo conceptual & \textbf{Deducción por descarte:} el motor de toda adivinanza. Tomás las pistas, tachás lo que no puede ser y te quedás con la única opción que sobrevive. Si solo queda una, esa es --- con certeza, no por suerte ---.\\
\end{tabularx}
\normalsize

\begin{drawbox}{milcTurquesa}{Cómo se resuelve un acertijo de grilla}
\textbf{1. Dibuja la grilla:} personas en las filas, objetos (o lugares) en las columnas. \\[1mm] \textbf{2. Lee cada pista} y marca con una \textbf{X} cada celda que esa pista descarta. \\[1mm] \textbf{3. Cuando en una fila queda una sola celda sin X,} ese es el objeto de esa persona: márcala con un círculo. \\[1mm] \textbf{4. Actualiza:} el objeto ya asignado se X en las otras filas. \\[1mm] \textbf{5. Repite} hasta que todas las filas tengan su objeto. Si no llegas, busca la pista que aún no usaste. En Bebras no gana quien calcula más rápido, sino quien \emph{descarta con más claridad}.
\end{drawbox}

\begin{softbox}{milcMostaza}{milcGris}{Errores comunes y su corrección}
\textbf{Leer el condicional al revés:} «SI A ENTONCES B» no significa «SI B ENTONCES A». De «vende mora» no se deduce que «vende lulo», aunque la regla diga lo contrario. \textbf{Olvidar actualizar la grilla:} cuando asignás un objeto a alguien, marca X en esa columna para los demás. \textbf{Saltar pistas:} cada pista descarta algo; si omitís una, puedes quedarte con dos opciones posibles en vez de una.
\end{softbox}

\begin{infoband}{milcTurquesa}
\textbf{En tu cuaderno (Actividad 2 · ANALIZA).} Título: «Actividad 2 --- Las 4 herramientas con mis palabras». \textbf{Formato:} tabla de 2 columnas (herramienta · ejemplo de mi barrio o mi casa). \textbf{Extensión:} las cuatro herramientas, cada una con un ejemplo distinto a los del texto. \textbf{Sabes que terminaste cuando} un compañero entendería tus ejemplos sin que se los expliques.
\end{infoband}

\puente{Ya distinguís las cuatro herramientas. Es hora de usarlas: vas a resolver un acertijo de grilla real y luego a juzgar si una conclusión ajena está bien deducida.}

\begin{titlebox}{milcMagenta}
Fase 3 · Praxis con pensamiento computacional
\end{titlebox}

\begin{softbox}{milcMagenta}{milcGris}{Construyendo el producto con los cuatro pilares}
Las herramientas no son adorno: se usan. Vas a resolver un \textbf{acertijo de grilla real} aplicando las cuatro a propósito, y luego a emitir un veredicto sobre si una deducción ajena es válida o no. Lee con calma --- la clave está en descartar con rigor, no en adivinar.
\end{softbox}

\small
\begin{tabularx}{\textwidth}{>{\bfseries\RaggedRight\arraybackslash}p{3.6cm}Y}
\rowcolor{milcMagenta!90}\color{white}Pilar & \color{white}\bfseries Aplicación al producto\\
Descomposición & \textbf{Condicional:} convertí cada pista en una regla «SI\ldots ENTONCES\ldots» y aplicala hacia adelante. Si la pista dice «Andrés no vive en Zaragoza», es «SI es Andrés ENTONCES no es Zaragoza»: marca X en esa celda.\\
Patrones & \textbf{Negación:} cada «no» en las pistas es tu aliado: te dice exactamente qué celda X marcar. Acumulá las X hasta que solo quede una opción sin tachar.\\
Abstracción & \textbf{Contrapositiva:} si la regla dice «todo el que baila, canta» y Pedro NO canta, entonces Pedro NO baila. Usa esta vuelta solo cuando tienes certeza de que la consecuencia es falsa.\\
Algoritmo de armado & \textbf{Descarte sistemático:} no adivinés ni saltes pasos. Grilla, pista por pista, X por X. Cuando quede una sola casilla limpia por fila, ya tienes la respuesta --- con certeza.\\
\end{tabularx}
\normalsize

\begin{drawbox}{milcMagenta}{El acertijo de los tres barrios}
\checkbox \textbf{Caso.} Andrés, Bruno y Camila viven cada uno en un barrio distinto de Cartago: El Prado, Santa Ana o Zaragoza. \\[1mm]
\checkbox \textbf{Pista 1:} Andrés no vive en Zaragoza. \\[1mm]
\checkbox \textbf{Pista 2:} Camila vive en El Prado. \\[1mm]
\checkbox \textbf{Pista 3:} el que vive en Zaragoza no es Camila. \\[2mm]
\checkbox Dibuja la grilla 3×3 y marca con X lo que cada pista descarta. \\[1mm]
\checkbox Escribe la asignación final: quién vive en qué barrio. \\[1mm]
\checkbox Explicá, en un párrafo, por qué cada persona no podía ir en los otros barrios.
\end{drawbox}

\begin{softbox}{milcMagenta}{milcGris}{Plantilla de trabajo}
\textbf{Solución paso a paso.} Pista 2: Camila = El Prado $\Rightarrow$ X en Camila/Santa Ana y Camila/Zaragoza; además El Prado se X para Andrés y Bruno. Pista 1: Andrés $\neq$ Zaragoza $\Rightarrow$ X en Andrés/Zaragoza. Solo queda Andrés/Santa Ana sin X $\Rightarrow$ Andrés = Santa Ana $\Rightarrow$ Bruno = Zaragoza. \textbf{Respuesta correcta:} Camila --- El Prado, Andrés --- Santa Ana, Bruno --- Zaragoza.
\end{softbox}

\begin{infoband}{milcMagenta}
\textbf{En tu cuaderno (Actividad 3 · EVALÚA).} Título: «Actividad 3 --- Veredicto sobre una deducción». \textbf{Formato:} veredicto (válida / inválida) + justificación + tu ejemplo corregido. \textbf{Extensión:} 6 a 8 renglones. \textbf{Sabes que terminaste cuando} tu veredicto dice «inválida», explicás que la regla no se puede leer al revés, y tu ejemplo corregido sí permite deducir algo con certeza.
\end{infoband}

\puente{Resolviste el acertijo y emitiste tu veredicto. Ahora deja tu evidencia: la grilla resuelta y la justificación escrita, para mostrar que sabes deducir con rigor.}

\begin{titlebox}{milcMostaza}
Producto de la sesión
\end{titlebox}
\textbf{Acertijo de los tres barrios resuelto con grilla + veredicto escrito sobre una deducción inválida}.
\begin{softbox}{milcMostaza}{milcGris}{Criterios de calidad}
(1) Tu grilla 3×3 muestra una X por cada descarte realizado, pista por pista, en el orden correcto. (2) Llegaste a la única asignación válida: Camila --- El Prado, Andrés --- Santa Ana, Bruno --- Zaragoza. (3) Tu párrafo explica, en orden, por qué cada persona no podía vivir en los otros barrios. (4) Tu veredicto dice que la conclusión «doña Rosa vende lulo porque vende mora» es \emph{inválida}, y justificás que la regla no se lee al revés. (5) Tu ejemplo corregido sí permite deducir algo con certeza (por ej.: «no vende mora» obliga a que tampoco venda lulo).
\end{softbox}

\puente{Terminaste el producto. Detente a mirar qué cambió en vos --- no la nota, sino tu manera de separar lo verdadero de lo falso.}

\begin{titlebox}{milcVino}
Fase 4 · Evaluación liberadora — cinco dimensiones
\end{titlebox}

\begin{softbox}{milcVino}{milcGris}{Sentido de la evaluación}
Evaluar no es solo revisar una nota. Es reconocer qué aprendí, qué cambió en mí, qué emociones manejé, cómo me relaciono con la ciudad y la comunidad, y qué le dejaré a quien viene detrás.
\end{softbox}

\textbf{Mi reflexión liberadora en cinco dimensiones.} Completa cada frase iniciada:

\small
\begin{tabularx}{\textwidth}{>{\bfseries\RaggedRight\arraybackslash}p{3.4cm}Y>{\RaggedRight\arraybackslash}p{4.6cm}}
\rowcolor{milcVino}\color{white}Dimensión & \color{white}\bfseries Mi respuesta (frame starter) & \color{white}\bfseries Algo que descubrí\\
1. Personal & Lo que más cambió en mí fue \linead{6.5cm} & \linead{4.2cm}\\
2. Emocional & La emoción más fuerte fue \linead{2.5cm} y la regulé \linead{3cm} & \linead{4.2cm}\\
3. Ciudadana & En democracia esto importa porque \linead{6cm} & \linead{4.2cm}\\
4. Local (Cartago) & En mi barrio sería útil para \linead{6.5cm} & \linead{4.2cm}\\
5. Intergeneracional & A mi abuela/abuelo le diría \linead{6.5cm} & \linead{4.2cm}\\
\end{tabularx}
\normalsize

\begin{infoband}{milcVino}
\textbf{Para la próxima sustentación.} Vuelve a esta tabla en seis meses. Verás que las respuestas cambian — no porque lo aprendido cambie, sino porque tú cambias. La evaluación liberadora es una conversación contigo a través del tiempo.
\end{infoband}


\puente{Cerramos pensando en grande: tres voces para entender qué significa, para vos y tu comunidad, aprender a razonar con certeza.}

\begin{titlebox}{milcGradeDark}
Triángulo de pensamiento · cierre filosófico
\end{titlebox}

\begin{softbox}{milcVino}{milcGris}{Dussel · lente del nosotros}
\textit{``El saber del pueblo también es ciencia.''}

La adivinanza del patio ya enseñaba a deducir con rigor --- descartar lo falso y quedarse con lo verdadero --- mucho antes de que alguien llamara eso «lógica formal». El pensamiento lógico no llegó de afuera: tu comunidad ya lo practicaba en los corredores y las noches de Cartago.

\textbf{¿Qué adivinanzas, acertijos o juegos de tu familia ya te enseñaron a deducir, mucho antes de que te lo explicaran en clase?}
\end{softbox}
\linead{\linewidth}\\[1mm]
\linead{\linewidth}

\begin{softbox}{milcMostaza}{milcGris}{Estoicismo · Marco Aurelio · cuidado interior}
\textit{``Todo lo que oímos es una opinión, no un hecho. Todo lo que vemos es una perspectiva, no la verdad.''}

Antes de creer una conclusión --- propia o ajena --- la lógica te pide detenerte y revisar: ¿de verdad se deduce de los hechos, o solo suena convincente? Esa pausa es el escudo estoico: no te apresurás a concluir, exigís la cadena completa de pistas.

\textbf{¿Me detengo a revisar si una conclusión mía o ajena de verdad se deduce de los hechos, o la acepto porque suena bien?}
\end{softbox}
\linead{\linewidth}\\[1mm]
\linead{\linewidth}

\begin{softbox}{milcTurquesa}{milcGris}{Floridi · lente de la infoesfera}
\textit{``La lógica es la gramática de la información: separa lo que se sigue de lo que solo lo parece.''}

En internet te llegan mil afirmaciones cada día. Saber distinguir un condicional bien aplicado de uno leído al revés es tu herramienta para separar información válida de ruido --- desde un acertijo Bebras hasta una noticia en redes.

\textbf{¿Con qué pistas y con qué lógica distinguís una afirmación verdadera de una falsa, en vez de creer la que más te gusta?}
\end{softbox}
\linead{\linewidth}\\[1mm]
\linead{\linewidth}

\begin{titlebox}{milcVino}
Compromiso final
\end{titlebox}

\small
\begin{tabularx}{\textwidth}{>{\RaggedRight\arraybackslash}p{3.0cm}YYY}
\rowcolor{milcVino}\color{white}\bfseries Criterio & \color{white}\bfseries Lo logré & \color{white}\bfseries En proceso & \color{white}\bfseries Necesito apoyo\\
Comprensión del tema & Explico ideas centrales con mis palabras. & Reconozco ideas pero falta precisión. & Necesito repasar conceptos base.\\
Pensamiento computacional & Apliqué los 4 pilares al diseñar mi producto. & Apliqué algunos pilares. & Necesito releer la fase 2.\\
Aplicación y producto & Mi producto cumple los criterios y se puede revisar. & Producto funcional con ajustes pendientes. & Aún en construcción.\\
Reflexión 5 dimensiones & Respondí cada dimensión con honestidad. & Respondí algunas con detalle. & Mis respuestas son muy generales.\\
Triángulo de pensamiento & Conecté las 3 voces con mi trabajo real. & Conecté al menos una. & No relaciono las voces con mi proceso.\\
\end{tabularx}
\normalsize

\textbf{Hoy entendí que:}\\[1mm]
\linead{\linewidth}\\[1mm]
\linead{\linewidth}

\textbf{Mi compromiso para la próxima sesión es:}\\[1mm]
\linead{\linewidth}\\[1mm]
\linead{\linewidth}

\begin{softbox}{milcMostaza}{milcGris}{Cita de despedida · Marco Aurelio}
\textit{``El obstáculo en el camino se vuelve el camino. Lo que estorba al camino se vuelve camino.''}\\[1mm]
Cada error de hoy, bien observado, se convierte en pieza del camino que pisarás mañana.
\end{softbox}

\small
\begin{tabularx}{\textwidth}{>{\bfseries}p{3.6cm}Y}
\rowcolor{milcGradeDark!12}Nombre completo: & \linead{10cm}\\
Fecha: & \linead{4cm} \quad Curso/grupo: \linead{4cm}\\
Mi firma: & \linead{9cm}\\
\end{tabularx}
\normalsize

\begin{infoband}{milcGradeDark}
\textbf{Para la próxima sesión.} Busca en casa o en el barrio una adivinanza que no conozcas --- preguntale a un familiar, busca en un libro de coplas o refranes colombianos --- y resuélvela por escrito usando la grilla o la cadena de descartes. Traé también un ejemplo propio de condicional que hayas visto en tu vida: una regla del salón, una norma de la casa, un letrero de la calle. Escribe la regla, un dato que la activa y lo que deducís.
\end{infoband}

\end{document}
